% !TeX root = ../main.tex

\documentclass
[	
	pdftex,	 	
	paper				= a4,				% paper size
	toc					= listof,			% list of figures and tables added to table of content
	toc					= bib,				% references section added to table of content
	titlepage			= true,	 			% set title in 'titlepage' environment
	numbers 			= noendperiod, 		% remove last point after enumeration
	draft 				= false,			
	twoside 			= false,	 		% modify page margin on both sides
	headinclude			= true,				% add header to text body
	footinclude			= false,			% add footer to text body
	fontsize 			= 11pt, 			% set document font size
]
{scrreprt}

% DO NOT EDIT THESE LINES!
\newif\iftikzexternal
\newif\ifgerman

%-------------------------------------------
%-- make your customized settings here! ----
%-------------------------------------------

%------------------------------------------------
%-- main settings that appear on the frontpage --
%------------------------------------------------

% document title and author
\newcommand{\documenttitle}{Entwicklung eines Gestaltungskonzepts für ein In-Car-Onboarding von Fahrzeug-Assistenzfunktionen zur Förderung einer sicheren und korrekten Nutzung}
\newcommand{\documentauthor}{\textbf{Jan Caspers}}
\newcommand{\studentid}{10019654}

% which type is your thesis? And which grade do you seek to achieve?
\newcommand{\thesistype}{MASTERARBEIT}
\newcommand{\thesisgrade}{Master of Science}

%\newcommand{\thesistype}{MASTERARBEIT}
%\newcommand{\thesisgrade}{Master of Science}

%\newcommand{\thesistype}{DISSERTATION}
%\newcommand{\thesisgrade}{Doctor rerum naturalium}

% in which subject/field do you graduate?
\newcommand{\thesisfield}{Mensch-Technik-Interaktion}

% where do you graduate?
\newcommand{\thesiswhere}{
    am \textbf{Fachbereich 1} \\
    der \textbf{Hochschule Ruhr West}
}

% where is your affiliation located?
\newcommand{\thesislocation}{Duisburg}

% who is the first examiner/reviewer and what is his/her affiliation?
\newcommand{\firstexaminer}{Prof. Dr. Stefan Geisler}
\newcommand{\firstexamineraffiliation}{Hochschule Ruhr West}

% who is the second examiner/reviewer and what is his/her affiliation?
\newcommand{\secondexaminer}{Prof. Dr. Stefan Becker}
\newcommand{\secondexamineraffiliation}{Hochschule Ruhr West}

%-------------------------------------------
%-- further template settings for latex ----
%-------------------------------------------

% date format for complete document
\usepackage{datetime}
\newdateformat{monthyeardate}{%
    \monthname[\THEMONTH] \THEYEAR
}%
\newdateformat{daymonthyeardate}{%
    \twodigit{\THEDAY}. \monthname[\THEMONTH] \THEYEAR
}

% place your images within images/ folder
\usepackage{graphicx}
\graphicspath{{images/}}

% comment in if your document is in German language
% otherwise, the document essentials are in English
\germantrue

% comment in if you want tikz to buffer the images as PDFs in the folder 'tikz/'
% this might significantly speed up LaTex compilation!
% Note: if commented-in, you also need to pass the argument '-shell-escape' to the LaTex compiler!
%\tikzexternaltrue

% add your custom packages here
\usepackage{makecell}
\usepackage{ifthen}
\usepackage{stackengine}
\usepackage{subcaption}
\usepackage{float}
\usepackage{longtable}
\usepackage{array}
\renewcommand\thesubfigure{\roman{subfigure}}

% set vertical distance after single paragraph
\parskip6pt

% set horizontal distance after single paragraph
\parindent6pt


%-------------------------------------------
%-- DO NOT MODIFY! Here starts the magic ---
%-------------------------------------------

%-------------------------------------------
%------------- Package section -------------
%-------------------------------------------

\usepackage[onehalfspacing]{setspace}
\usepackage[utf8]{inputenc}
\usepackage{listings}
\usepackage{tabularx}
\usepackage{bbm}
\usepackage{abstract}
\usepackage{multirow}
\usepackage{amsmath}
\usepackage{amssymb}
\usepackage{dsfont}
\usepackage[usenames, dvipsnames]{xcolor}
\usepackage{framed}
\usepackage{booktabs}
\usepackage{acronym}
\usepackage{caption}
\usepackage{hyperref}
\usepackage{enumitem}
\usepackage[intoc, noprefix]{nomencl}
\makenomenclature

\ifgerman
    \usepackage[ngerman]{babel}
    \renewcommand{\nomname}{Nomenklatur}
\fi


% last but not least, initialize tikz library
\usepackage{tikz}
\usepackage{overpic}

% if 'tikzexternal' is true, enable tikz buffer mode
\iftikzexternal
    \usetikzlibrary{external}
    \tikzexternalize[prefix=tikz/]
    
    % significantly increase the memory buffer to prevent memory overflow
    \tikzset{external/system call={pdflatex \tikzexternalcheckshellescape -extra-mem-top=10000000 -extra-mem-bot=10000000 -halt-on-error -interaction=batchmode -jobname "\image" "\texsource"}}
\fi

\usepackage{pgfplots}
\pgfplotsset{compat=1.15}


%-------------------------------------------
%--------- Page properties section ---------
%-------------------------------------------

% set margin of single page
\usepackage[
    paper=a4paper,
    left=25mm,
    right=20mm,
    top=30mm,
    bottom=40mm
]{geometry}

% set width of head separator line
\usepackage[headsepline=0.5pt]{scrlayer-scrpage}
\usepackage{scrhack}

% set page style to scrheadings and clear default properties to overwrite with custom properties
\pagestyle{scrheadings}
\clearpairofpagestyles

% automatically update title in headline based on the current chapter
\automark{chapter}

% set head line caption on the right-hand side of the page
\ohead{\rightmark}

% overwrite font familiy for head line
\setkomafont{pageheadfoot}{\normalfont\sffamily\bfseries\small}
\renewcommand\chapterheadstartvskip{\vspace*{-1cm}}

% page enumerations - center page numbering and set font family
\cfoot[\pagemark]{\pagemark}
\setkomafont{pagenumber}{\normalfont\sffamily}

% set the font familiy which overwrites some defaults
\renewcommand{\rmdefault}{stix} %This overwrites the roman default
\renewcommand{\sfdefault}{stix} %This overwrites the sans serif default
\renewcommand{\ttdefault}{cmtt} %This overwrites the monospaced (“typewriter”) default
\renewcommand*\familydefault{\rmdefault}

% switch environment so that '@' is a 'standard' letter
\makeatletter

% redefine the spacing after a subsubsection call
\let\origsubsubsection\subsubsection
\renewcommand\subsubsection{\@ifstar{\starsubsubsection}{\nostarsubsubsection}}

\newcommand\nostarsubsubsection[1]
{\subsubsectionprelude\origsubsubsection{#1}\subsubsectionpostlude}

\newcommand\starsubsubsection[1]
{\subsubsectionprelude\origsubsubsection*{#1}\subsubsectionpostlude}

% spacing before subsubsection
\newcommand\subsubsectionprelude{% 
    \vspace{0em}
}

% spacing after subsubsection
\newcommand\subsubsectionpostlude{% 
    \vspace{-1em}
}

%-------------------------------------------
%------ lstlisting properties section ------
%-------------------------------------------

% use arabic numbering for (code) lstlisting
\AtBeginDocument	
{																
    \renewcommand*{\thelstlisting}{\arabic{lstlisting}}			
}

% switch back so that '@' is a special character in the following
\makeatother													

% set some further properties for lstlisting
\renewcommand{\lstlistingname}{Source Code} 
\lstset{
    captionpos 			= b,                   				% place caption at the bottom of the listing
    numbers 			= left,            					% place line count on the left
    numbersep 			= 5pt,           					% 5pt distance to source code
    numberstyle 		= \tiny,       						% font size 'tiny' for enumeration
    breaklines 			= true,         					% break lines if required
    breakautoindent 	= true,    							% auto indentation after line break
    tabsize 			= 2,               					% tab size 2
    frame 				= shadowbox,						% shadow box
    rulesepcolor 		= \color{Gray},						% color of shadow
    framexleftmargin 	= 5mm,								% set margin of left frame
    basicstyle 			= \ttfamily\footnotesize,			% \footnotesize \small \scriptsize
    keywordstyle 		= \color{Red}\bfseries,				% color of key words
    commentstyle 		= \color{ForestGreen}\bfseries, 	% color of comments
    stringstyle 		= \color{NavyBlue},					% color of strings
    showspaces			= false,        					% do not show white spaces
    showstringspaces	= false   							% also do not show white spaces in strings
}


%-------------------------------------------
%---------- PDF meta information -----------
%-------------------------------------------

\hypersetup{					
    pdftitle			= 	{\documenttitle},   			% document title
    pdfauthor			= 	{\documentauthor},     			% document author
    pdfsubject			=	{\documenttitle},   			% subect of the document, same as the title
    bookmarksnumbered	= 	true, 							% enable enumeration for booktabs
    colorlinks			= 	true,       					% color links?
    linkcolor			= 	Black,          				% color of links
    citecolor			= 	Black,        					% color of citations
    urlcolor			= 	Black,    						% color of URLs
}

%-------------------------------------------
%--------- Miscellaneous settings ----------
%-------------------------------------------

% use this command to buffer the actual page count of Roman page numbering
% so that the page count can be reused within the reference section
\ExplSyntaxOn%
\DeclareExpandableDocumentCommand{\arabicnumeral}{m}
{
    \int_from_roman:n { #1 }
}
\ExplSyntaxOff

%-- tabularx special column type --%
\newcolumntype{Y}{>{\centering\arraybackslash}X}

% distance between equations within same environment
\setlength{\jot}{10pt} 
